\section{Data and protocol}
\label{sec:protocol}

\paragraph{Data.} The PKDD'99 financial data set \citep{berka2000guide},
from the CTU relational repository \citep{motl2015ctu}, holds the
transactions of \R{ho/data/accounts_total} accounts of a Czech bank from
\R{berka/first_year} to \R{berka/last_year}. Activity is sparse and booked:
the median held-out account has \R{ho/data/median_records_per_month}
records a month, with no pending state.
- The end-of-day balance is the running sum of amounts from the opening
  balance. The bank's balance field is used only to validate it.
- Two adapter decisions were committed before any development result: an
  every-day holiday calendar, and the exclusion of accounts whose running
  sum drifts from the reported balances by more than \R{c/max-drift}
  crowns.
- Of the \R{ho/data/accounts_in_split} held-out accounts,
  \R{ho/data/accounts_excluded_balance_inconsistent} are excluded for
  drift, and \R{ho/data/accounts_with_eligible_origin} have eligible
  origins (\R{ho/data/eligible_origins} origins).
- Amounts are rescaled per origin by $k = \R{c/reference-scale}/s$, where
  $s$ is the account's median monthly gross flow.

Appendix~\ref{app:protocol} gives these details, including a recount of the
drift rule. We also use two synthetic generators and one UK account
(MoneyData, Appendix~\ref{app:moneydata}).

\paragraph{Origins and methods.} Origins fall every \R{c/origin-step} days
from day \R{c/warmup} of an account's history. Each forecast uses only the
records dated on or before its origin. A test that cut the input at
\R{au/prefix/origins} random development origins reproduced all
\R{au/prefix/truncated_equals_full} forecasts exactly.
- \textbf{V2} is the primary model, frozen.
- \textbf{B-EQ} is today's balance plus the account's own past $h$-day
  changes over the last year.
- \textbf{B-SA+}, the plan's reference, is the detected schedule plus the
  empirical distribution of its own past errors.
- \textbf{B-GLOBAL} is gradient-boosted quantile regression
  \citep{ke2017lightgbm} pooled across accounts.

Four further per-account baselines and the scoring details are in
Appendix~\ref{app:protocol}.

\paragraph{Scores.}
- Coverage of central intervals, pooled.
- CRPS, and the probability integral transform (PIT).
- CRPS skill of $m$ against a reference, given two ways over the forecasts
  both made:
  - the \emph{account mean}
    $\operatorname{mean}_a(1-\sum_a \mathrm{CRPS}_m/\sum_a
    \mathrm{CRPS}_{\text{ref}})$, the planned estimand;
  - the \emph{pooled} ratio over all origins.
- Intervals are \R{c/ci-level}\% intervals from \R{c/bootstrap}
  resamples of whole accounts.
- All methods share one origin set: \R{ho/var/V2/30/forecasts} held-out
  forecasts in \R{ho/var/V2/30/accounts} accounts per horizon.

\paragraph{The analysis plan, the revision and the freeze.} The plan fixed
falsifiable claims before any real-data run. Its two central claims:
- \textbf{A2:} on held-out data, every central interval within
  $\pm\R{c/a2-tol}$ points of nominal at \hA\ and \hB\ days and
  $\pm\R{c/a2-tol-90}$ at \hC;
- \textbf{B1:} CRPS skill against \emph{each} baseline positive at \hB\ and
  \hC\ days, with an interval excluding zero.

On development data, the original model V0 failed A2: its \lC\% interval
covered \R{dev/var/V0/90/cov90}\% at \hC\ days. A revision rule written
before any revision code chose among four candidates:
- \textbf{V1:} conformal recalibration \citep{gibbs2021aci};
- \textbf{V2:} compound-Poisson discretionary spending;
- \textbf{V3:} pooled priors;
- \textbf{V4:} V2 and V3 combined, recalibrated as V1.

The rule required \lC\% coverage within $\pm\R{c/select-tol}$ points at
every horizon, then chose the lowest CRPS. It selected V2, which was
frozen.

The held-out split runs only through a command that refuses an
uncommitted working tree, logs each run, and refuses a second run without
a logged reason. The main held-out evaluation ran once.

\paragraph{What is exploratory.} Everything reported after that run is
labelled exploratory, with no confirmatory status:
- a second, logged run of the planned ablations, which run~1 omitted
  through a harness gap;
- post-hoc summaries of the stored forecasts (paired comparisons, pooled
  skills, one-sided misses, per-cell verdicts);
- the analyses behind \S\ref{sec:explanation} and \S\ref{sec:overdraft}.
  They were specified in writing (the repository's
  \texttt{paper/REVISION.md}), developed on the development split, and
  then run on held-out accounts in two further logged runs.

\paragraph{Timeline.}
\label{sec:timeline}
Git commit times are set by the author's machine; GitHub's record of when
it first received each commit is not.
- The freeze was on GitHub at \R{tl/70a7bca/pushed}, before held-out run~1
  began at \R{tl/run1/time}.
- The held-out results were on GitHub at \R{tl/cdf070d/pushed}, before the
  ablation run began at \R{tl/run2/time}.
- The plan was first pushed at \R{tl/8c40faf/pushed}, together with the
  first development results. Its precedence over them rests on its
  author-controlled commit time, \R{tl/8c40faf/time}.

The full timeline and every deviation are in Appendix~\ref{app:protocol}.
